\section{Linked Flows and Worker Outcomes}\label{app:flows}

\subsection{Link construction and risk sets}

Person links follow official IPUMS/Census matching guidance and use the available longitudinal weight for each eligible link.  Adjacent-month and twelve-month endpoints are distinct samples.  A twelve-month link is not the sum of twelve observed monthly transitions and may conceal intervening changes.  Table~\ref{tab:app_flow_retention} reports fixed-support retention ranges across exposure quintiles rather than collapsing different periods and risk sets into one number.

Employment exit begins with an employed person in an occupation on the analysis support.  Reported occupation change requires employment at both endpoints and a valid occupation at each.  Entry-destination allocation begins with a nonemployed origin, ends with observed employment, and asks whether the destination falls in Q5 rather than Q1.  It therefore conditions on entry and is neither an employment-finding probability nor an employer hiring rate.  Every crosswalk descendant of one source record shares its link and resampling weight.

\begin{table}[!htbp]
\centering
\caption{Official-weight link retention across exposure quintiles}\label{tab:app_flow_retention}
\begin{tabular}{llrr}
\toprule
Horizon & Origin age & Preperiod range & Postperiod range \\
\midrule
Adjacent month & Ages 22--25 & 87.6--88.9\% & 85.3--86.9\% \\
Adjacent month & Ages 26--65 & 90.5--92.4\% & 87.6--90.5\% \\
Twelve month & Ages 22--25 & 50.1--54.0\% & 48.7--52.8\% \\
Twelve month & Ages 26--65 & 68.8--73.3\% & 63.9--70.2\% \\
\bottomrule
\end{tabular}
\tabnote{Each entry is the minimum--maximum across fixed beta quintiles among employed origins on the 468-occupation analysis support.  Numerators require a validated endpoint and a positive horizon-specific longitudinal weight; denominators use origin \texttt{WTFINL}.  Rotation ineligibility and a target month outside the extract are recorded separately from eligible origins without a validated endpoint.}
\end{table}


Retained and eligible-unlinked origins differ observably.  Among postperiod
older Q5 origins at twelve months, retained records average age 44.69 and have
a 59.25-percent BA-or-higher share, versus 41.50 and 53.35 percent among
eligible-unlinked records.  Among postperiod young Q2 origins, the 35-or-more
hours share is 60.83 percent retained versus 66.24 percent unlinked.  These
comparisons use origin \texttt{WTFINL}; they diagnose observable selection but
do not correct selection or sign bias from unobserved outcomes.

\subsection{Results and uncertainty}

At the adjacent-month horizon, the young-relative Q5--Q1 estimate is 0.132 for employment exit, 0.003 for reported occupation change, and $-0.070$ for entry-destination allocation.  Their respective 95-percent intervals are $[-0.042,0.307]$, $[-0.118,0.123]$, and $[-0.250,0.110]$.  At twelve months, the three estimates are 0.123, $-0.018$, and $-0.052$.  Their intervals are $[-0.111,0.357]$, $[-0.097,0.061]$, and $[-0.251,0.147]$.  All include zero.  Person- and household-cluster sensitivities, descriptive transition rates, effective risk counts, and named occupation influence appear in the machine-readable flow directory.

Usual-hours changes are estimated within linked employed persons.  The adjacent estimate is $-0.119$ hours per week with interval $[-0.295,0.057]$.  Weekly earnings use the eligible outgoing-rotation sample and its earnings weight; the conditional log-earnings coefficient is $-0.0245$ with interval $[-0.1024,0.0534]$.  The earnings extract covers 74 months through March 2023 and therefore contains only three postperiod months.  Both outcomes condition on employment.  Unemployment duration is described only among persons observed unemployed, and no occupation is imputed to never-employed or out-of-labor-force respondents.  These restrictions prevent the outcomes from being misreported as population employment effects.

\begin{table}[!htbp]
\centering
\caption{Employment-exit components in linked CPS records}\label{tab:app_exit_components}
\resizebox{\textwidth}{!}{%
\begin{tabular}{llrrr}
\toprule
Endpoint & Component & Raw events & Coefficient & 95\% interval \\
\midrule
Adjacent month & Unemployment entry & 32,488 & $0.055$ & $[-0.272,0.381]$ \\
Adjacent month & Labor-force exit & 74,480 & $0.135$ & $[-0.064,0.333]$ \\
Twelve month & Unemployment at endpoint & 27,330 & $0.246$ & $[-0.295,0.787]$ \\
Twelve month & Out of labor force at endpoint & 78,773 & $0.078$ & $[-0.181,0.337]$ \\
\bottomrule
\end{tabular}}
\tabnote{Rows use official longitudinal weights and the same young-relative Q5--Q1 log-rate-ratio contrast as the core flow table.  Counts are repeated link events rather than unique persons.  Unemployment and out-of-labor-force events exhaust observed employment exits, but their nonlinear coefficients do not add.  Every interval includes zero.}
\end{table}


\subsection{Selection sensitivity and missing-outcome bounds}

Table~\ref{tab:app_flow_selection} compares three descriptive linear
probability contrasts.  Official link weights and origin-weight complete cases
use the retained population directly.  Observable poststratification instead
reweights within fixed observed strata and therefore requires conditional
missing-at-random.  Factors are neither trimmed nor capped; their maximum is
102.7.  Similarity across those three columns cannot rule out selection on
unobservables.

\begin{table}[!htbp]
\centering
\caption{Linked-flow selection sensitivities}\label{tab:app_flow_selection}
\scriptsize
\setlength{\tabcolsep}{3pt}
\resizebox{\textwidth}{!}{%
\begin{tabular}{llrrrr}
\toprule
Horizon & Margin & Official link & Origin-weight complete case & Observable poststratification & Linear accounting interval \\
\midrule
Adjacent & Employment exit & 0.0088 & 0.0088 & 0.0086 & $[-0.441,0.463]$ \\
Adjacent & Unemployment entry & 0.0037 & 0.0039 & 0.0038 & $[-0.445,0.459]$ \\
Adjacent & Labor-force exit & 0.0051 & 0.0050 & 0.0048 & $[-0.444,0.460]$ \\
Twelve month & Employment exit & 0.0166 & 0.0167 & 0.0210 & $[-1.570,1.575]$ \\
Twelve month & Unemployment entry & 0.0119 & 0.0128 & 0.0130 & $[-1.573,1.572]$ \\
Twelve month & Labor-force exit & 0.0047 & 0.0039 & 0.0080 & $[-1.576,1.569]$ \\
\bottomrule
\end{tabular}}
\tabnote{The first three numerical columns are descriptive linear-probability Q5--Q1, young--older, post--pre differences, not nonlinear regression coefficients.  Observable poststratification assumes outcome missingness is independent of linkage within fixed observed strata.  Accounting intervals impose no conditional-MAR assumption or Lee trimming and may exceed the probability-contrast parameter space because they bound the linear regression contrast jointly over source records.}
\end{table}


The last column assigns one common binary missing outcome to each eligible
source record and all of its bridge descendants.  Its linear interval is sharp
under that source-record coupling for the named margin.  No conditional-MAR
assumption or Lee trimming is used.  The analogous log intervals are only
conservative marginal outer intervals because bridge coupling need not preserve
joint sharpness.  These accounting intervals concern missing outcomes, whereas
Main Table 5's confidence intervals concern sampling uncertainty.

\subsection{Entry denominators and annual timing}

Table~\ref{tab:app_entry_probabilities} places every observed destination on
one retained-nonemployment denominator.  The destination probabilities sum to
one after including people who remain nonemployed and employed destinations
outside the exposure support.  The conditional Q1--Q5 allocation used in the
regression divides only by observed supported entries and remains a different
object; neither is employer hiring.

\begin{table}[!htbp]
\centering
\caption{Employment-entry destinations on a common risk set}\label{tab:app_entry_probabilities}
\scriptsize
\setlength{\tabcolsep}{2pt}
\resizebox{\textwidth}{!}{%
\begin{tabular}{lllrrrrrrrr}
\toprule
Horizon & Age & Period & Remain nonemployed & Q1 & Q2 & Q3 & Q4 & Q5 & Outside support & Any entry \\
\midrule
Adjacent & 26--65 & Pre & .9154 & .0221 & .0180 & .0117 & .0087 & .0113 & .0128 & .0846 \\
Adjacent & 26--65 & Post & .9179 & .0210 & .0175 & .0115 & .0091 & .0105 & .0125 & .0821 \\
Adjacent & 22--25 & Pre & .8491 & .0361 & .0398 & .0191 & .0129 & .0211 & .0219 & .1509 \\
Adjacent & 22--25 & Post & .8626 & .0361 & .0352 & .0162 & .0117 & .0193 & .0190 & .1374 \\
Twelve month & 26--65 & Pre & .8169 & .0439 & .0396 & .0252 & .0213 & .0263 & .0268 & .1831 \\
Twelve month & 26--65 & Post & .8276 & .0399 & .0385 & .0230 & .0215 & .0246 & .0250 & .1724 \\
Twelve month & 22--25 & Pre & .6408 & .0772 & .0885 & .0492 & .0390 & .0589 & .0465 & .3592 \\
Twelve month & 22--25 & Post & .6611 & .0695 & .0843 & .0555 & .0351 & .0530 & .0415 & .3389 \\
\bottomrule
\end{tabular}}
\tabnote{The denominator is all retained nonemployed origins, weighted with the horizon-specific official link weight.  Q1--Q5 are employed destinations on the retained exposure support.  Outside support records an employed endpoint with a valid occupation outside that support.  Any entry equals one minus remaining nonemployment and reconciles to the sum of all observed employed destinations.}
\end{table}


The annual exclusion rule drops origins whose twelve-month window straddles
January 2023 and is reported with the other primary margins in
main Table 5.  The exposure-duration sensitivity retains those origins
and codes the fraction of the twelve destination months at or after onset.  Age
and exposure remain defined at the origin.  Table~\ref{tab:app_annual_timing}
reports only this additional sensitivity, avoiding a second finite-draw
interval for the primary target.  No paired interval comparing the two rules
is available, so their difference is not described as detected or equivalent.
October 2025 is absent rather than interpolated, and annual endpoints remain
distinct from sums of monthly transitions.

\begin{table}[!htbp]
\centering
\caption{Twelve-month exposure-duration sensitivity}\label{tab:app_annual_timing}
\begin{tabular}{lr}
\toprule
Margin & Exposure-duration coding \\
\midrule
Employment exit & $0.158\;[-0.070,0.387]$ \\
Unemployment entry & $0.186\;[-0.337,0.708]$ \\
Labor-force exit & $0.144\;[-0.084,0.372]$ \\
Reported occupational outflow & $-0.008\;[-0.082,0.067]$ \\
\bottomrule
\end{tabular}
\tabnote{Entries are Q5--Q1, young--older onset coefficients under exposure-duration coding, which codes the share of the twelve destination months at or after onset (102 months).  The baseline exclusion rule, which drops origins whose twelve-month link straddles January 2023, is the primary specification and is reported in main Table 5 (89 months).  This column is not paired with that table.  Intervals use 9,999 occupation wild-score draws.}
\end{table}


\subsection{Why no stock--flow calibration is reported}

A compatible accounting identity would require entries from nonemployment, exits, switches among all occupations, aging into and out of ages 22--25, corresponding movements for the older denominator, job-to-job moves, and population/sample-composition residuals.  The CPS endpoint links do not measure all of these components on a common employer-match risk set, and the BCC hiring and separation objects are not identical to CPS employment transitions.  A single implied stock response would therefore be assumption driven.  The revision reports the measurable margins and the precise comparability gap instead of manufacturing a calibration.
